\begin{center}\begin{tikzpicture}[font=\small,>=stealth,line join=round]
\draw[gray!40,rounded corners=3pt] (0,0) rectangle (14.2,8.3);\node at (7.1,7.700000000000001) {例3：内层虽反号，外层平方把两路合并};
\begin{scope}[shift={(4.1,2.161111111111111)},x=1.875cm,y=0.5111111111111111cm]\draw[->,gray] (-1.6,0)--(1.6,0) node[right]{$x$};\draw[->,gray] (0,-1)--(0,8) node[above]{$y$};\draw[blue,thick,solid,domain=-1.3:1.3,samples=180] plot(\x,{4*(\x)^2});\fill[orange!85!black] (-1,4) circle(2.2pt);\fill[orange!85!black] (1,4) circle(2.2pt);\end{scope}\node[anchor=west,align=left,text width=5.0cm] at (8.4,4.0) {$F=(2x)^2$\\[10pt]$g(1)=2,\ g(-1)=-2$\\[10pt]平方以后都变成 $4$\\[10pt]内奇外偶，整体偶};
\node[text width=12.6cm,align=center] at (7.1,.65) {不能把“内层为奇”误记成“整体必奇”；还需要检查外层。};
\end{tikzpicture}\end{center}
